\scanpage{33}
structure \inserted{de vecteur} complexe sur $\operatorname{DR}(M)$ ; mais on voit que la bigraduation déduite par transport de structure de celle de Hodge ne \uline{dépend pas de $\pi$} (donc, si on remplace $\pi(\lambda)$ par $\overline{\pi(\lambda)}$, elle n'est \uline{pas} remplacée par la symétrique !).

\inserted{Néanmoins, si on regarde, pour un motif donné $M$, le motif $\unsure{{}^{c}M_{\underline c}}$ ($\underline c$ l'aut. non trivial de $K$) sur $\operatorname{DR}({}^{c}M_{\underline c})\xrightarrow{\sim}\operatorname{DR}(M)_{\underline c}$, les structures \unsure{rationnelles}, \inserted{(et filtrations)} et les $\unsure{\sigma}$, \unsure{Ainsi|Alors} \deleted{la filtration} composant avec $\sigma$, on trouve un \uline{$K$-isom} $\operatorname{DR}(M_{\underline c})\simeq\operatorname{DR}(M)$, conservant la $\mathbf Q$-structure, et \uuline{renversant} \inserted{les} \unsure{filtrations}.}

\deleted{\unsure{Supp.}} Soit maintenant $K_0=\mathbf R$, et considérons motifs \inserted{$M$} \uline{sur} $\mathbf R$. Choisissons une clôture alg. $K$ de $K_0=\mathbf R$. Alors $\operatorname{DR}(M)_K=\operatorname{DR}(M_K)$ est muni de \uline{deux façons} d'une $\mathbf R$-structure : du fait qu'il est déduit de $\operatorname{DR}(M)$ par $\mathbf R\unsure{-|\longrightarrow}K$, et puis : la construction donnée plus haut. Soit $\sigma$ l'autom de conjugaison complexe défini plus haut, (N.B. $\sigma$ \uline{ne respecte pas} la \uline{filtration}) $\tau$ celle définie par $\operatorname{DR}(M)_K\simeq\operatorname{DR}(M_K)$ (donc $\unsure{\sigma|\tau}$ respecte la filtration) - $\sigma$ et $\tau$ sont \uline{semilinéaires}. Alors $\sigma$ et $\tau$ commutent, en d'autres termes $\operatorname{DR}(M_K)^\sigma$ $(\simeq H^*(X(K),\mathbf R)$ si $M=\operatorname{Mot}^*(X))$ est stable par $\tau$ : en effet, $\tau$ n'est autre que l'autom de $H^*(X(K),\mathbf R)$ défini par l'autom $\tau_X$ de $X(K)$ ! Alors $\sigma\tau=\tau\sigma=\unsure{f_{\mathbf R}|\mathfrak f_{\mathbf R}}$ est l'autom de $H^*(X(K),K)\simeq\operatorname{DR}(\operatorname{Mot}^*(X))$

\margin{$X(K)$}
\mspage{32}
